\documentclass[12pt,leqno]{amsart}
\usepackage{pifont}
\usepackage{stmaryrd}
%\usepackage{amsfonts}
%\usepackage{amscd}
\usepackage{dsfont}
\usepackage{color}
\usepackage{mathrsfs}
\usepackage{amsmath}
\usepackage{amssymb}
\usepackage[hidelinks]{hyperref}
\usepackage{bookmark}
\usepackage[bottom]{footmisc}
\usepackage{verbatim}
\usepackage{extarrows}
\usepackage{mathtools}
\usepackage{orcidlink}

\oddsidemargin -.1in \evensidemargin -.1in \textwidth 6.5in
\textheight 8.2in
\linespread{1.3}

\numberwithin{equation}{section}
\newtheorem{thm}{Theorem}[section]
\newtheorem{lem}[thm]{Lemma}
\newtheorem{ex}[thm]{Example}
\newtheorem{prop}[thm]{Proposition}
\newtheorem{cor}[thm]{Corollary}
\newtheorem{defi}[thm]{Definition}
\newtheorem{rmk}[thm]{Remark}
\newtheorem{notation}[thm]{Notation}

\newcommand{\End}{\operatorname{End}}
\newcommand{\Hom}{\operatorname{Hom}}
\newcommand{\Ext}{\operatorname{Ext}}
\newcommand{\Soc}{\operatorname{Soc}}
\newcommand{\rad}{\operatorname{rad}}
\newcommand{\ch}{\operatorname{ch}}
\newcommand{\one}{\mathbf{1}}
\newcommand{\id}{\operatorname{id}}
\newcommand{\wt}{\operatorname{wt}}
\newcommand{\xqz}[1]{\lfloor #1 \rfloor}
\newcommand{\Res}{\operatorname{Res}}
\newcommand{\Aut}{\operatorname{Aut}}
\newcommand{\Vir}{\mathfrak{Vir}}
\newcommand{\TL}{\mathcal{TL}}
\newcommand{\VirMod}{\mathcal L}
\newcommand{\boxt}{\boxtimes}
\newcommand{\KL}{KL}
\newcommand{\Oc}{\mathcal O}
\newcommand{\Kac}{\mathcal K}
\renewcommand{\mod}{\operatorname{mod}}

\def\Z{\mathbb{Z}}
\def\N{\mathbb{N}}
\def\A{\mathbb{A}}
\def\C{\mathbb{C}}
\def\R{\mathbb{R}}
\def\Q{\mathbb{Q}}
\def\F{\mathbb{F}}

\allowdisplaybreaks

\begin{document}

\title{Quantum Drinfeld--Sokolov Reduction and the McRae--Yang Tensor Functors}

\author{Shun Xu~\orcidlink{0009-0006-8080-8107}}

\address{School of Mathematical Sciences, Anhui University, Anhui, Hefei, 230601, China}

\email{shunxu@ahu.edu.cn}

\subjclass[2020]{Primary 17B69, 17B68; Secondary 18M15, 81R10}

\keywords{Kazhdan--Lusztig category, Virasoro algebra, quantum Drinfeld--Sokolov reduction, Kac module, logarithmic module, projective object}

\begin{abstract}
Let $p,q\geq2$ be relatively prime and set $k=-2+p/q$ and
$c_{p,q}=1-6(p-q)^2/(pq)$.  McRae and Yang constructed a right exact
braided tensor functor
\[
 F_{p,q}:\KL^k(\mathfrak{sl}_2)\longrightarrow\Oc_{c_{p,q}}
\]
sending the standard Weyl module $V_2$ to the Virasoro Kac module
$\Kac_{2,1}$, and conjectured that it is exact and naturally isomorphic to
quantum Drinfeld--Sokolov reduction.  We prove this conjecture.  We first
determine the images of every indecomposable projective under $F_{p,q}$
without assuming exactness of $F_{p,q}$.  The proof combines the projective
tensor recursions in $\KL^k(\mathfrak{sl}_2)$, exact fusion by $\Kac_{2,1}$,
and generalized conformal-weight blocks.  Projective faithfulness and the
non-semisimple affine twists then imply that these images are logarithmic.
Using Nakano's logarithmic Virasoro extension theorem and his computation of
the relevant $\Ext^1$-groups, we identify the projective images and compute
their Hom spaces, proving full faithfulness on projectives.  Quantum
Drinfeld--Sokolov reduction has the same projective images.  Compatibility
with the non-standard affine twist fixes the scalars needed to construct a
natural isomorphism on projectives, and uniqueness of right exact extension
then gives
\[
 F_{p,q}\cong H^0_{DS,+}\big|_{\KL^k(\mathfrak{sl}_2)}.
\]
In particular, $F_{p,q}$ is exact.  The analogous assertion holds after
interchanging $p$ and $q$.
\end{abstract}

\maketitle

\section{Introduction}

Let $p,q\geq2$ be relatively prime integers and put
\begin{equation}\label{eq:parameters}
 k=-2+\frac pq,
 \qquad
 c_{p,q}=1-\frac{6(p-q)^2}{pq}.
\end{equation}
Throughout the proof, $(p,q)$ is an \emph{ordered} pair: every statement
through Proposition~\ref{prop:natural-on-projectives} is proved for an
arbitrary ordered pair of relatively prime integers at least two.  Thus the
same proved statement may subsequently be applied, without a new
``transposed'' argument, to the ordered pair $(q,p)$.
The category $\KL^k(\mathfrak{sl}_2)$ of finite-length grading-restricted
generalized modules for the universal affine vertex algebra
$V^k(\mathfrak{sl}_2)$ is a non-semisimple braided tensor category.  McRae
and Yang determined its projective objects and identified its projective
subcategory with the tilting category for quantum $\mathfrak{sl}_2$ at the
root of unity $e^{\pi iq/p}$ \cite{McRaeYang}.  Their universal property
produces a right exact braided tensor functor
\begin{equation}\label{eq:F-intro}
 F_{p,q}:\KL^k(\mathfrak{sl}_2)\longrightarrow\Oc_{c_{p,q}},
 \qquad F_{p,q}(V_2)\cong\Kac_{2,1},
\end{equation}
where $\Oc_{c_{p,q}}$ is the braided tensor category of finite-length
Virasoro modules studied in \cite{McRaeSopin}.  McRae and Yang conjectured
that $F_{p,q}$ is exact and agrees naturally with quantum
Drinfeld--Sokolov reduction \cite[Conjecture~7.16]{McRaeYang}.

The obstacle is subtle.  The source projective $P_{np+r}$ admits a Weyl
filtration
\[
 0\longrightarrow V_{np-r}\longrightarrow P_{np+r}
 \longrightarrow V_{np+r}\longrightarrow0,
\]
but applying a merely right exact functor does not preserve the injection on
the left.  Consequently, projective images cannot be read off from this
sequence.  We instead recover them from tensor recursions.  This is the first
main step of the proof.

The second main step is a Virasoro extension calculation.  Once the
projective image is known to have the expected Kac filtration, preservation
of the non-standard twist makes it logarithmic.  Nakano's Theorem~5.11 and
Proposition~5.12 in \cite{Nakano} then give precisely
the two facts needed here: an appropriate quotient of the logarithmic
extension is indecomposable, and the relevant extension space is
one-dimensional.  These statements permit a direct computation of all Hom
spaces between projective images.  Importantly, our argument never assumes
that a general Nakano module $P(\tau)$ has a prescribed socle or is
self-contragredient.  Nakano explicitly observes that the general socle
statement is not supplied by his construction
\cite[Figure~5.1 and the discussion following it]{Nakano}.

Our main result is the following.

\begin{thm}\label{thm:main}
Let $p,q\geq2$ be relatively prime and let $k$ and $c_{p,q}$ be as in
\eqref{eq:parameters}.  The McRae--Yang functor
\[
 F_{p,q}:\KL^k(\mathfrak{sl}_2)\longrightarrow\Oc_{c_{p,q}}
\]
is exact, and there is a natural isomorphism of $\C$-linear functors
\[
 F_{p,q}\cong H^0_{DS,+}\big|_{\KL^k(\mathfrak{sl}_2)}.
\]
Likewise,
\[
 F_{q,p}:\KL^{-2+q/p}(\mathfrak{sl}_2)\longrightarrow\Oc_{c_{p,q}}
\]
is exact and is naturally isomorphic to the corresponding principal quantum
Drinfeld--Sokolov reduction.  Thus Conjecture~7.16 of McRae--Yang holds.
\end{thm}

The proof is organized as follows.  Section~\ref{sec:prelim} fixes notation
and records the source projectives, one-row Virasoro Kac modules, and the two
Nakano extension inputs.  Section~\ref{sec:DS} identifies the reductions of
Weyl and simple modules and proves compatibility of reduction with the
non-standard affine twist.  Section~\ref{sec:images} determines all
projective images of $F_{p,q}$ without using exactness.  In
Section~\ref{sec:fullness}, logarithmicity and Nakano's results yield the
target Hom algebra and hence full faithfulness on projectives.  Finally,
Section~\ref{sec:comparison} compares the projective images with
Drinfeld--Sokolov reduction, constructs a natural isomorphism by a
twist-locking argument, and extends it uniquely to the entire abelian
category.

\section{Preliminaries}\label{sec:prelim}

\subsection{Affine modules and projectives}

Write $V_r$ for the Weyl module induced from the $r$-dimensional simple
$\mathfrak{sl}_2$-module and $L_r$ for its simple quotient.  We use
$P_r\twoheadrightarrow L_r$ for the projective cover.  The following facts
are proved in \cite[Sections~2--4]{McRaeYang}.  The category
$\KL^k(\mathfrak{sl}_2)$ has enough projectives.  For $1\leq r\leq p-1$,
and also whenever $p\mid r$, one has $P_r=V_r$; in particular
$P_{np}=V_{np}=L_{np}$ is simple-projective.  If $n\geq1$ and
$1\leq r\leq p-1$, then
\begin{equation}\label{eq:source-projective-filtration}
 0\longrightarrow V_{np-r}\longrightarrow P_{np+r}
 \longrightarrow V_{np+r}\longrightarrow0
\end{equation}
is non-split, and $P_{np+r}$ has composition factors
\begin{equation}\label{eq:source-composition-factors}
 [P_{np+r}]=2[L_{np+r}]+[L_{np-r}]+[L_{(n+2)p-r}].
\end{equation}
The tensor products $V_2\boxt P_j$ satisfy the projective recursions of
\cite[Theorem~4.8]{McRaeYang}; these will be used in the proof of
Proposition~\ref{prop:projective-images}.

The category is $\Z/2\Z$-graded: if $M$ lies in degree $\bar\epsilon$, then
the $h_0$-eigenvalues on $M$ lie in $\epsilon+2\Z$.  In addition to the
standard twist, McRae and Yang use
\begin{equation}\label{eq:minus-twist}
 \theta^-_M=(-1)^\epsilon e^{2\pi iL^{\mathrm{aff}}_0}.
\end{equation}
Their functor $F_{p,q}$ preserves this twist
\cite[Theorem~7.15]{McRaeYang}.  For $n\geq1$ and $1\leq r\leq p-1$,
Theorem~4.13 of \cite{McRaeYang} says that $P_{np+r}$ is logarithmic, while
the projective Hom algebra described in Sections~4 and 6 there gives
\[
 \End(P_{np+r})=\C\id\oplus\C N_{np+r},
 \qquad N_{np+r}^2=0,
\]
with $N_{np+r}\neq0$ spanning the radical.  Consequently
\begin{equation}\label{eq:source-twist-nilpotent}
 \theta^-_{P_{np+r}}
 =\lambda_{np+r}(\id+\gamma_{np+r}N_{np+r}),
 \qquad \gamma_{np+r}\neq0.
\end{equation}
The proof of Theorem~4.13 of \cite{McRaeYang} identifies the nonzero radical
direction with the nilpotent part of $L^{\mathrm{aff}}_0$.  Hence its exponential has
nonzero coefficient $\gamma_{np+r}$ in the displayed basis.

\subsection{Virasoro Kac modules}

Let $\Kac_{r,s}$ denote the Virasoro Kac module of central charge
$c_{p,q}$, and abbreviate
\[
 K_j:=\Kac_{j,1},
 \qquad
 S_j:=L(c_{p,q},h_{j,1}),
 \qquad
 h_j:=h_{j,1}.
\]
Thus
\begin{equation}\label{eq:weights}
 h_j=\frac{(qj-p)^2-(p-q)^2}{4pq}.
\end{equation}
The category $\Oc_{c_{p,q}}$ is an abelian category of finite-length
grading-restricted generalized Virasoro modules.  Its Hom spaces are
finite-dimensional: a map between simple highest-weight modules is
determined by its restriction to the one-dimensional lowest-weight space,
and the assertion for arbitrary finite-length objects follows by induction
on the two lengths.  Since
idempotents split in an abelian category, $\Oc_{c_{p,q}}$ is therefore
Krull--Schmidt.  In this paper, a generalized module is called
\emph{logarithmic} precisely when $L_0$ is not semisimple; no additional
condition on its Loewy shape is included in this term.  The Kac structure results of
\cite[Section~2.3]{McRaeSopin} imply that $K_{np}$ is simple and, for
$n\geq0$ and $1\leq r\leq p-1$,
\begin{equation}\label{eq:Kac-length-two}
 0\longrightarrow S_{(n+2)p-r}\longrightarrow K_{np+r}
 \longrightarrow S_{np+r}\longrightarrow0
\end{equation}
is non-split.  We put $K_0=0$.

Fusion by the rigid module $K_2$ is exact, and McRae--Sopin prove
\begin{equation}\label{eq:K2-fusion}
 0\longrightarrow K_{j-1}\longrightarrow K_2\boxt K_j
 \longrightarrow K_{j+1}\longrightarrow0;
\end{equation}
this sequence splits if and only if $p\nmid j$
\cite[Theorem~6.5]{McRaeSopin}.

\subsection{Nakano's extension input}

We recall only the portion of Nakano's notation needed below.  Set
\[
 p_+=\min\{p,q\},\qquad p_-=\max\{p,q\}.
\]
Write $\mathscr L_{p_+,p_-}=L_{c_{p_+,p_-}}\text{-mod}$ for Nakano's
category in Definition~5.1 of \cite{Nakano}.  It is the full subcategory of
generalized Virasoro modules having finite socle length and all simple
factors in the Felder weight set $H_{p_+,p_-}$.  For the finite-length
objects used here this subcategory is extension-closed: if
$0\to M'\to M\to M''\to0$ has finite-length end terms
$M',M''\in\mathscr L_{p_+,p_-}$, then $M$ has finite length and its
Jordan--H\"older factors, with multiplicity, are those of $M'$ and $M''$.
Thus $M$ has finite socle length and all its simple factors again have
weights in $H_{p_+,p_-}$.
The simples occurring in $\Oc_{c_{p,q}}$ have lowest weights in this Felder
set, so this observation applies to every module and every short exact
sequence used below.  Thus all Ext groups in this subsection are Yoneda Ext
groups in exactly Nakano's category, not Ext groups imported from a larger
module category.  Nakano's Definition~3.14 also calls a module
logarithmic exactly when $L_0$ acts non-semisimply, agreeing with our
convention above.

Let $\tau=(\alpha_1,\alpha_2,\alpha_3)$ be a triple in Nakano's set
$\mathcal T_{p_+,p_-}$: the corresponding Fock modules are consecutive terms of a
single Felder complex and
\[
 h_{\alpha_1}\leq h_{\alpha_2}<h_{\alpha_3}.
\]
When the first inequality is strict, write
\[
 A(\tau)=L(h_{\alpha_1},h_{\alpha_2}),
 \qquad B(\tau)=L(h_{\alpha_2},h_{\alpha_3})
\]
for the unique non-split length-two modules, and let $K(\tau)$ be the
corresponding indecomposable extension of $B(\tau)$ by
$L(h_{\alpha_1})$ as in \cite[Definitions~5.9--5.10]{Nakano}.
The asserted uniqueness follows from the one-dimensional simple--simple
Ext groups in \cite[Proposition~5.5]{Nakano}; hence it is uniqueness in
$\mathscr L_{p_+,p_-}$, not merely uniqueness among modules with a chosen
presentation.

We shall use the following two results, quoted exactly in the form needed.

\begin{thm}[Nakano]\label{thm:Nakano-input}
Let $\tau\in\mathcal T_{p_+,p_-}\setminus\mathcal T^0_{p_+,p_-}$.  If $E$ is a
logarithmic module representing a class in
\[
 \Ext^1_{\mathscr L_{p_+,p_-}}\bigl(B(\tau),A(\tau)\bigr),
\]
then $E/L(h_{\alpha_2})$ is indecomposable.  Moreover, for every
$\tau\in\mathcal T_{p_+,p_-}$,
\[
 \Ext^1_{\mathscr L_{p_+,p_-}}
 \bigl(K(\tau),L(h_{\alpha_2})\bigr)\cong\C.
\]
Consequently the middle term $P(\tau)$ of the unique nonzero class is
well-defined up to isomorphism.
\end{thm}

\begin{proof}
The first assertion is \cite[Theorem~5.11]{Nakano}; the second is
\cite[Proposition~5.12 and Definition~5.13]{Nakano}.
\end{proof}

\begin{rmk}\label{rmk:no-socle}
No statement about the socle or contragredient self-duality of a general
$P(\tau)$ will be used.  Indeed, Nakano notes that the dotted socle arrow in
the general diagram cannot be determined from the arguments preceding
Figure~5.1 of \cite{Nakano}.
\end{rmk}

We now match the triples used below with Nakano's labels.  Recall from
\cite[(2.3)--(2.4)]{Nakano} that
\begin{equation}\label{eq:Nakano-weight-shift}
 h_{u,v;m}=h_{u-mp_+,v}=h_{u,v+mp_-}.
\end{equation}

For $n\geq1$ and $1\leq r\leq p-1$, set
\begin{equation}\label{eq:abc}
 a=np-r,\qquad b=np+r,\qquad c=(n+2)p-r.
\end{equation}

\begin{lem}[Felder-label matching]\label{lem:Felder-labels}
There is an explicitly labelled triple
\(\tau_{n,r}\in\mathcal T_{p_+,p_-}\setminus\mathcal T^0_{p_+,p_-}\) whose ordered
conformal weights are \((h_a,h_b,h_c)\).  More precisely, if $p<q$, then
\begin{equation}\label{eq:tau-p-less-q}
 \tau_{n,r}=
 \bigl(\alpha_{r,q-1;n-1},
       \alpha_{p-r,q-1;n},
       \alpha_{r,q-1;n+1}\bigr),
\end{equation}
whereas if $p>q$, then
\begin{equation}\label{eq:tau-p-greater-q}
 \tau_{n,r}=
 \bigl(\alpha_{q-1,r;-n+1},
       \alpha_{q-1,p-r;-n},
       \alpha_{q-1,r;-n-1}\bigr).
\end{equation}
In either case,
\begin{equation}\label{eq:tau-nr}
 (h_{\alpha_1},h_{\alpha_2},h_{\alpha_3})=(h_a,h_b,h_c)
\end{equation}
and
\begin{equation}\label{eq:A-B-tau}
 A(\tau_{n,r})\cong K_a,
 \qquad B(\tau_{n,r})\cong K_b.
\end{equation}
\end{lem}

\begin{proof}
Suppose first that $p<q$, so $(p_+,p_-)=(p,q)$.  Part~(1) of Nakano's
Felder complex \cite[Proposition~2.7]{Nakano}, with the second label fixed
at $q-1$, contains the consecutive segment
\[
 F_{r,q-1;n-1}\longrightarrow F_{p-r,q-1;n}
 \longrightarrow F_{r,q-1;n+1}.
\]
Using \eqref{eq:Nakano-weight-shift}, or directly the momentum formula in
\cite[(2.1)--(2.4)]{Nakano}, gives respectively
\[
 h_{r,q-1;n-1}=h_{np-r,1},\quad
 h_{p-r,q-1;n}=h_{np+r,1},\quad
 h_{r,q-1;n+1}=h_{(n+2)p-r,1}.
\]
This proves the label and weight assertions in \eqref{eq:tau-p-less-q}.

If $p>q$, then $(p_+,p_-)=(q,p)$.  Part~(2) of the same proposition,
with first label $q-1$, contains the consecutive segment
\[
 F_{q-1,r;-n+1}\longrightarrow F_{q-1,p-r;-n}
 \longrightarrow F_{q-1,r;-n-1}.
\]
The symmetry $h^{(p,q)}_{j,1}=h^{(q,p)}_{1,j}$ and
\eqref{eq:Nakano-weight-shift} give the same ordered weights
$(h_a,h_b,h_c)$.  Thus \eqref{eq:tau-p-greater-q} also satisfies the
adjacency requirement in Nakano's Definition~5.6.

Finally, \eqref{eq:weight-difference} (or direct substitution into
\eqref{eq:weights}) gives
\[
 h_b-h_a=r(nq-1)>0,
 \qquad
 h_c-h_b=(p-r)((n+1)q-1)>0.
\]
Hence the first inequality in Nakano's definition is strict, so the triple
does not lie in $\mathcal T^0_{p_+,p_-}$.  The identifications
\eqref{eq:A-B-tau} now follow from the two non-split Kac sequences
\eqref{eq:Kac-length-two} and uniqueness of the corresponding length-two
extensions.
\end{proof}

\section{Drinfeld--Sokolov reduction}\label{sec:DS}

Let
\[
 H:=H^0_{DS,+}\big|_{\KL^k(\mathfrak{sl}_2)}.
\]

We first justify both the domain and the exactness of this restriction.  This
is needed because objects of $\KL^k(\mathfrak{sl}_2)$ need not have
semisimple affine conformal Hamiltonian.

\begin{prop}\label{prop:H-exact}
For every $M\in\KL^k(\mathfrak{sl}_2)$,
\[
 H^i_{DS,+}(M)=0\qquad(i\neq0).
\]
Consequently $H$ is exact, and its values belong to $\Oc_{c_{p,q}}$.
\end{prop}

\begin{proof}
Each simple object $L_r$ is an affine highest-weight module and hence an
object of the affine BGG category to which the BRST vanishing theorem
applies.  Theorem~9.1.3 of \cite{ArakawaW} gives vanishing outside degree
zero under condition (383), and Remark~9.1.5(ii) there states that in type
$A_1$ no condition on the block is required.  The resulting unconditional
exactness of $H^0_{DS,+}$ on affine category $O$ in type $A_1$ is also stated
directly in \cite[Section~4.4, discussion preceding formula~(4.20)]
{ArakawaFrenkel}.  Thus
\[
 H^i_{DS,+}(L_r)=0\qquad(i\neq0)
\]
for every $r\geq1$.

The BRST construction is a functorial cochain complex on the smooth affine
modules under consideration.  Therefore every short exact sequence in
$\KL^k(\mathfrak{sl}_2)$ gives a long exact sequence in BRST cohomology.
Induction on Jordan--H\"older length, using a short exact sequence with
simple quotient, propagates the preceding vanishing from simples to every
finite-length generalized object.  The same long exact sequence then
reduces to a short exact sequence in degree zero, so $H$ is exact.  Notice
that this argument does not require the generalized object itself to have
semisimple $L^{\mathrm{aff}}_0$.

Finally, Theorem~9.1.4 and Remark~9.1.5(ii) of \cite{ArakawaW} identify
$H(L_r)$ with a nonzero simple Virasoro highest-weight module.  Hence, by
exactness and another induction on length, $H(M)$ is a finite-length
grading-restricted generalized Virasoro module whose simple factors are the
modules $H(L_r)$.  Their lowest weights are in the extended Kac table by
the calculation in Proposition~\ref{prop:H-simple} below.  Thus
$H(M)\in\Oc_{c_{p,q}}$.
\end{proof}

\begin{prop}\label{prop:H-Weyl}
For every $r\geq1$,
\[
 H(V_r)\cong K_r.
\]
\end{prop}

\begin{proof}
Let $\widehat\lambda=(r-1)\omega+k\Lambda_0$ and let
$M(\widehat\lambda)$ be the affine Verma module.  The finite-dimensional
$r$-dimensional $\mathfrak{sl}_2$-module is a quotient of its finite Verma
module, and induction gives a canonical surjection
\[
 M(\widehat\lambda)\twoheadrightarrow V_r.
\]
Both modules in this surjection belong to the affine BGG category $O$.
The exactness statement in Theorem~9.1.3 of \cite{ArakawaW}, with the block
hypothesis removed in type $A_1$ by Remark~9.1.5(ii) there, therefore gives
a surjection; equivalently, one may use the direct exactness statement in
\cite[Section~4.4, discussion preceding formula~(4.20)]{ArakawaFrenkel}:
\[
 H^0_{DS,+}(M(\widehat\lambda))\twoheadrightarrow H(V_r).
\]
For $\mathfrak{sl}_2$, quantum reduction sends the affine Verma module to
the Virasoro Verma module with transformed highest weight; see
\cite[Section~4.4, especially the discussion preceding formula~(4.20)]
{ArakawaFrenkel} and the underlying theorem in
\cite{ArakawaSuperconformal}.  Thus
$H^0_{DS,+}(M(\widehat\lambda))\cong
\mathcal V(c_{p,q},h_r)$.  This proves, rather than assumes, that $H(V_r)$
is cyclic.  Formula~\eqref{eq:H-character} below has coefficient one in
weight $h_r$, so the image of the Verma highest-weight vector is nonzero;
thus $H(V_r)$ is generated specifically by its lowest-weight BRST class.

It remains to identify the quotient.  Take $\lambda=(r-1)\omega$ and
$\check\mu=0$ in the rank-one specialization of
\cite[formula~(4.19)]{ArakawaFrenkel}; the cohomological concentration needed
to read the Euler characteristic as the character of $H^0$ is
\cite[Theorem~2.1]{ArakawaFrenkel}.  This gives
\begin{equation}\label{eq:H-character}
 \ch H(V_r)=
 \frac{z^{h_r}(1-z^r)}{\prod_{m\geq1}(1-z^m)}.
\end{equation}
Hence the kernel has no vector below level $r$ and is one-dimensional at
level $r$.  A nonzero vector of that level in the kernel is therefore
singular.  A singular vector at any fixed level of a Virasoro Verma module
is unique up to scalar \cite[Proposition~2.2.1(3)]{CreutzigEtAl}; moreover,
$K_r$ is the quotient by the submodule generated by the level-$r$ singular
vector \cite[Remark~2.2]{McRaeSopin}.  Thus the original surjection factors through
a surjection $K_r\twoheadrightarrow H(V_r)$.  The character of $K_r$ is
the right-hand side of \eqref{eq:H-character}, so this surjection is an
isomorphism.
\end{proof}

\begin{prop}\label{prop:H-simple}
For every $r\geq1$,
\[
 H(L_r)\cong S_r.
\]
\end{prop}

\begin{proof}
For the $+$ reduction, Arakawa proves that the reduction of a simple
highest-weight module is the corresponding nonzero simple $W$-module under
the block condition in Theorem~9.1.4 of \cite{ArakawaW}.  In type $A_1$ no
block condition is required; this is stated explicitly in
\cite[Remark~9.1.5(ii)]{ArakawaW}.  Thus $H(L_r)$ is nonzero and simple.
Put $m=r-1$.  In the rank-one normalization used for the $+$ reduction, the
transformed highest weight in Theorem~9.1.4 has conformal weight
\[
 \Delta_{DS}(m)
 =\frac{m(m+2)}{4(k+2)}-\frac m2
 =\frac{q(r^2-1)-2p(r-1)}{4p},
\]
where $k+2=p/q$.  On the other hand, direct expansion of
\eqref{eq:weights} gives
\[
 h_r=h_{r,1}^{(p,q)}
 =\frac{(qr-p)^2-(p-q)^2}{4pq}
 =\frac{q(r^2-1)-2p(r-1)}{4p}.
\]
Hence $\Delta_{DS}(r-1)=h_r$, and the nonzero simple reduction is
$H(L_r)\cong S_r$.
\end{proof}

\begin{cor}\label{cor:H-faithful}
The functor $H$ is faithful.
\end{cor}

\begin{proof}
Let $0\neq f:M\to N$.  The nonzero finite-length module $\operatorname{Im}f$
has a simple subquotient $L_r$, and Proposition~\ref{prop:H-simple} shows
that its image under the exact functor $H$ is nonzero.  Consequently
\[
 \operatorname{Im}H(f)=H(\operatorname{Im}f)\neq0.
\]
\end{proof}

We next record the twist convention carefully.

\begin{prop}\label{prop:H-twist}
If $M\in\KL^k_{\bar\epsilon}(\mathfrak{sl}_2)$, then
\[
 H(\theta^-_M)=\theta_{H(M)},
\]
where $\theta^-_M$ is defined in \eqref{eq:minus-twist} and the twist on the
right is $e^{2\pi iL^{DS}_0}$.
\end{prop}

\begin{proof}
Write the module BRST complex as
\[
 C(M)=M\otimes\mathcal F_{\mathrm{gh}}.
\]
Let $L^{\mathrm f}_0$ be the standard conformal Hamiltonian on the charged
fermion factor.  The diagonal action of
$\bar\rho^\vee=h/2$ on $C(M)$ is
\begin{equation}\label{eq:diagonal-rho-action}
 \bar\rho^\vee_0=\frac12h_0+J^{\mathrm{gh}}_0,
\end{equation}
where $J^{\mathrm{gh}}_0$ is the ghost Cartan charge: on a ghost generator
associated with a root $\beta$, it acts by
$\beta(h/2)\in\{1,-1\}$.  Thus $J^{\mathrm{gh}}_0$ has integral spectrum.

The old conformal Hamiltonian on $C(M)$ is
$L^{\mathrm{aff}}_0+L^{\mathrm f}_0$.  Arakawa changes the degree operator
to $D_{\mathrm{new}}=D+\bar\rho^\vee$ and proves that
$-D_{\mathrm{new}}$ is the new Hamiltonian commuting with the $+$ BRST
differential \cite[(169)--(173), Section~4.6]{ArakawaW}.  Substituting the
old Hamiltonian and \eqref{eq:diagonal-rho-action} gives, on the module
complex,
\begin{equation}\label{eq:BRST-Hamiltonian}
 L^{DS}_0=L^{\mathrm{aff}}_0-\frac12h_0+
 \bigl(L^{\mathrm f}_0-J^{\mathrm{gh}}_0\bigr).
\end{equation}
We use the charged-fermion vacuum on which both $L^{\mathrm f}_0$ and
$J^{\mathrm{gh}}_0$ vanish.  Fermion creation modes change both eigenvalues
integrally, so $L^{\mathrm f}_0-J^{\mathrm{gh}}_0$ has integral spectrum.
The $h_0$-eigenvalues on $M$ lie in
$\epsilon+2\Z$, and therefore
\[
 e^{-\pi ih_0}=(-1)^\epsilon\id_M.
\]
All exponentials here are well defined on the generalized modules in this
paper.  Indeed, each generalized $L^{\mathrm{aff}}_0$-eigenspace is
finite-dimensional, and on it the exponential is the scalar exponential
times the finite exponential polynomial of the locally nilpotent part;
$h_0$, $L^{\mathrm f}_0$, and $J^{\mathrm{gh}}_0$ act semisimply on the
weight-space and ghost factors.  Exponentiating
\eqref{eq:BRST-Hamiltonian} therefore gives on the BRST complex
\[
 e^{2\pi iL^{DS}_0}=(-1)^\epsilon e^{2\pi iL^{\mathrm{aff}}_0}.
\]
Both sides commute with the BRST differential, so the equality descends to
cohomology.  This is exactly the claimed compatibility.  The charged
fermion conventions and the $+$ BRST operator are those of
\cite[(154), (166)--(168)]{ArakawaW}; the corresponding module complexes are
the ones used in \cite{ArakawaVanishing}.
\end{proof}

\section{Images of projective modules}\label{sec:images}

From now on, write $F=F_{p,q}$.  This section does not use exactness of
$F$.

\begin{lem}\label{lem:F-small-Weyl}
For $1\leq r\leq p$,
\[
 F(V_r)\cong K_r.
\]
\end{lem}

\begin{proof}
The assertions for $r=1,2$ follow from the tensor unit and the definition
of $F$.  Suppose $2\leq r\leq p-1$ and the result holds through $r$.
The affine sequence
\[
 0\longrightarrow V_{r-1}\longrightarrow V_2\boxt V_r
 \longrightarrow V_{r+1}\longrightarrow0
\]
splits because $p\nmid r$ \cite[Theorem~2.23]{McRaeYang}.  Hence
\[
 K_{r-1}\oplus F(V_{r+1})
 \cong K_2\boxt K_r
 \cong K_{r-1}\oplus K_{r+1},
\]
where the second isomorphism is the split case of
\eqref{eq:K2-fusion}.  Krull--Schmidt cancellation gives the conclusion.
\end{proof}

\begin{prop}\label{prop:F-faithful-projectives-large}
If $p\geq3$, the restriction of $F$ to the projective subcategory is
faithful.
\end{prop}

\begin{proof}
Put $\zeta=e^{\pi iq/p}$.  Theorem~6.6 of \cite{McRaeYang} gives a
monoidal equivalence
\begin{equation}\label{eq:tilting-projective-equivalence}
 \Phi:\mathcal T_\zeta\xrightarrow{\sim}\mathcal P^k,
 \qquad \Phi(X)=V_2,
 \qquad \Phi(T_\lambda)=P_{\lambda+1},
\end{equation}
where $\mathcal T_\zeta$ is the quantum $\mathfrak{sl}_2$ tilting category;
the last correspondence is Proposition~6.4 there.  Let
$I=\ker(F|_{\mathcal P^k})$ and pull it back along $\Phi$ to a tensor ideal
$I'$ of $\mathcal T_\zeta$.

Assume $I'\neq0$.  Taking a nonzero matrix component, choose
$0\neq f:T\to U$ in $I'$ with $T,U$ indecomposable.  Every indecomposable
tilting object is a direct summand of a tensor power of the standard object
$X$; see the construction of $\mathcal T_\zeta$ in
\cite[Section~6.1]{McRaeYang}.  Choose split embeddings and projections
\[
 T\xrightarrow{i_T}X^{\otimes m}\xrightarrow{p_T}T,
 \qquad
 U\xrightarrow{i_U}X^{\otimes n}\xrightarrow{p_U}U.
\]
Then $i_Ufp_T\neq0$: otherwise multiplication on the left by $p_U$ and on
the right by $i_T$ would give $f=0$.  Quantum Schur--Weyl duality in type
$A_1$ says that the pivotal functor
\begin{equation}\label{eq:TL-fullness}
 \TL(-\zeta-\zeta^{-1})\longrightarrow\mathcal T_\zeta,
 \qquad 1\longmapsto X,
\end{equation}
is full on all Hom spaces between tensor powers of $X$.  More precisely,
the functor induces an equivalence after additive idempotent completion
\cite[Proposition~2.13 and Remark~2.15]{SuttonTubbenhauerWedrichZhu}.
Hence $i_Ufp_T$ is the image of a nonzero Temperley--Lieb morphism, and the
pullback of $I'$ along
\eqref{eq:TL-fullness} is a nonzero tensor ideal.  It is proper, since the
tensor unit is not killed by $F$.

Goodman--Wenzl's ideal theorem now applies to this pullback.  By
\cite[Theorem~3.3]{GoodmanWenzl}, the pullback ideal is the negligible
ideal, and by \cite[Proposition~2.1]{GoodmanWenzl}, that ideal is generated
by the first critical Jones--Wenzl idempotent.  To match their parameter convention, set
$\tau=-\zeta$.  Then the loop parameter is
\[
 \tau+\tau^{-1}=-\zeta-\zeta^{-1},
 \qquad \tau^2=\zeta^2,
\]
and $\operatorname{ord}(\tau^2)=p$ because $(p,q)=1$.  Consequently their
first critical index is $p$ and the generator is the idempotent
$p_{[p-1]}$.  In the additive Karoubi envelope, the object
$(p-1,p_{[p-1]})$ corresponds to the indecomposable tilting module
$T_{p-1}$, and $p_{[p-1]}$ is the identity morphism of this Karoubi object.
Since a tensor ideal is closed under pre- and postcomposition, the
containment of $p_{[p-1]}$ therefore gives
\(\id_{T_{p-1}}\in I'\).  Now
\eqref{eq:tilting-projective-equivalence} implies \(\id_{P_p}\in I\).
This is impossible because
\[
 F(P_p)=F(V_p)\cong K_p\neq0
\]
by Lemma~\ref{lem:F-small-Weyl}.  Thus the kernel is zero.
\end{proof}

\begin{rmk}\label{rmk:TL-p-ge-3}
The hypothesis $p\geq3$ is essential for the preceding proof.  When
$p=2$, the loop parameter is zero; neither the stated range of the
Goodman--Wenzl ideal theorem nor the quantum-characteristic convention of
\cite{SuttonTubbenhauerWedrichZhu} covers that boundary case.  We give a
separate proof after determining the projective images.
\end{rmk}

We isolate the conformal-weight calculation driving the induction.

\begin{lem}\label{lem:weight-arithmetic}
For arbitrary integers $u,v$,
\begin{equation}\label{eq:weight-difference}
 h_u-h_v=\frac{(u-v)(q(u+v)-2p)}{4p}.
\end{equation}
In particular, for $n\geq1$ and $1\leq r\leq p-1$,
\begin{align}
 h_{np+r-1}-h_{np-r+1}&=(r-1)(nq-1),\label{eq:parallel-one}\\
 h_{np+r+1}-h_{np-r-1}&=(r+1)(nq-1),\label{eq:parallel-two}\\
 h_{np+r-1}-h_{np-r-1}&=r(nq-1)-\frac{rq}{p},\label{eq:cross-one}\\
 h_{np+r+1}-h_{np-r+1}&=r(nq-1)+\frac{rq}{p}.\label{eq:cross-two}
\end{align}
The first two differences are integral and the last two are not.  At a
$p$-wall the analogous two pairings are
\[
 \{K_{(n-1)p},K_{(n+1)p}\},\qquad
 \{K_{(n-1)p+2},K_{(n+1)p-2}\},
\]
and crossed differences have nonintegral parts
$\pm(p-1)q/p$.  If $p=2$, the corresponding nonintegral parts are
$\pm q/2$.
\end{lem}

\begin{proof}
Formula \eqref{eq:weight-difference} follows by subtracting
\eqref{eq:weights}.  The displayed identities follow by substitution.
Because $(p,q)=1$, neither $rq/p$ for $1\leq r\leq p-1$ nor
$(p-1)q/p$ is integral.  If $p=2$, then $q$ is odd.  The Virasoro modes
change generalized conformal weights by integers, so these calculations
give the asserted $\C/\Z$-blocks.
\end{proof}

\begin{lem}[Exact block projection]\label{lem:block-projection}
Let $M$ be a finite-length generalized Virasoro module.  For
$\xi\in\C/\Z$, let $M^{[\xi]}$ be the sum of the generalized $L_0$-weight
spaces whose weights belong to $\xi$.  Then
\[
 M=\bigoplus_{\xi\in\C/\Z}M^{[\xi]}
\]
with only finitely many nonzero summands, each $M^{[\xi]}$ is a Virasoro
submodule, and $M\mapsto M^{[\xi]}$ is exact.
\end{lem}

\begin{proof}
Since $L_m$ changes generalized $L_0$-weight by the integer $-m$, the stated
sum is by Virasoro submodules.  Finite length leaves only finitely many
congruence classes.  A Virasoro homomorphism preserves generalized
$L_0$-weights, so kernels and images decompose class by class.  Projecting a
short exact sequence to one direct summand is therefore again exact.
\end{proof}

\begin{prop}[Projective-image theorem]\label{prop:projective-images}
For $n\geq1$ and $1\leq r\leq p-1$,
\begin{equation}\label{eq:F-projective-filtration}
 0\longrightarrow K_{np-r}\longrightarrow F(P_{np+r})
 \longrightarrow K_{np+r}\longrightarrow0.
\end{equation}
Moreover,
\begin{equation}\label{eq:F-wall-projective}
 F(P_{np})\cong K_{np}
 \qquad(n\geq1).
\end{equation}
\end{prop}

\begin{proof}
We give the simultaneous induction in detail, including the point that
avoids assuming left exactness of $F$.  The initial wall object is
$Q_p=F(P_p)=K_p$ by Lemma~\ref{lem:F-small-Weyl}.  The source recursion
$V_2\boxt P_{np}\cong P_{np+1}$ and the nonsplit case of
\eqref{eq:K2-fusion} show that whenever $Q_{np}=K_{np}$,
\[
 0\longrightarrow K_{np-1}\longrightarrow Q_{np+1}
 \longrightarrow K_{np+1}\longrightarrow0.
\]
This starts the induction away from a wall.

Suppose next that the filtration is known for $Q_{np+r}$ with
$1\leq r\leq p-2$.  Tensor it with $K_2=F(V_2)$.  This operation is exact
by rigidity of $K_2$, and both outer Kac modules split under
\eqref{eq:K2-fusion}.  The four resulting factors form precisely the two
$\C/\Z$-blocks
\[
 \{K_{np-r+1},K_{np+r-1}\},
 \qquad
 \{K_{np-r-1},K_{np+r+1}\}
\]
by Lemma~\ref{lem:weight-arithmetic}.  On the other hand, the source
projective recursion gives
\begin{equation}\label{eq:interior-projective-recursion}
 K_2\boxt Q_{np+r}
 \cong Q_{np+r-1}\oplus Q_{np+r+1}
 \qquad(2\leq r\leq p-2),
\end{equation}
with the $r=1$ modification
\begin{equation}\label{eq:first-projective-recursion}
 K_2\boxt Q_{np+1}\cong2Q_{np}\oplus Q_{np+2}.
\end{equation}
By Lemma~\ref{lem:block-projection}, projection to either block is exact.
For $r\geq2$, the first block has precisely the two Kac factors
$K_{np-r+1}$ and $K_{np+r-1}$, which are exactly the factors of the already
known direct summand $Q_{np+r-1}$.  Since both modules have finite length,
the quotient of that block by this summand has Jordan--H\"older length zero
and is therefore zero.  Hence the entire summand $Q_{np+r+1}$ lies in the
second block, where exact block projection supplies its asserted short
exact sequence.  For $r=1$, the same argument uses the first block
$2K_{np}=2Q_{np}$ in \eqref{eq:first-projective-recursion}; its finite-length
complement is again zero.  This proves all non-wall cases in the $n$th
strip without any hidden left exactness.

At its far edge, apply $K_2\boxt-$ to $Q_{(n+1)p-1}$.  The source recursion is
\begin{equation}\label{eq:wall-projective-recursion}
 K_2\boxt Q_{(n+1)p-1}
 \cong Q_{(n-1)p}\oplus Q_{(n+1)p-2}\oplus Q_{(n+1)p},
\end{equation}
where $Q_0=0$.  Lemma~\ref{lem:weight-arithmetic} separates the Kac factors
into the two blocks
\[
 \{K_{(n-1)p},K_{(n+1)p}\},
 \qquad
 \{K_{(n-1)p+2},K_{(n+1)p-2}\}.
\]
The second block and the already known $Q_{(n+1)p-2}$ have the same two
Kac factors, so the same Jordan--H\"older length argument identifies them.
The first block is $Q_{(n-1)p}\oplus Q_{(n+1)p}$ and has factors
$K_{(n-1)p}$ and $K_{(n+1)p}$.  Removing the known simple summand
$Q_{(n-1)p}=K_{(n-1)p}$ leaves an object of length one with factor
$K_{(n+1)p}$; it is therefore isomorphic to that simple module.  Thus
$Q_{(n+1)p}=K_{(n+1)p}$.  This establishes the next wall and completes the
induction for $p\geq3$.

If $p=2$, start with $Q_2=K_2$.  The recursion
$K_2\boxt Q_{2n}=Q_{2n+1}$ gives
$0\to K_{2n-1}\to Q_{2n+1}\to K_{2n+1}\to0$.  Tensoring this sequence by
$K_2$ and using
\[
 K_2\boxt Q_{2n+1}
 \cong Q_{2(n-1)}\oplus2Q_{2n}\oplus Q_{2(n+1)}
\]
separates the factors into the blocks
$\{K_{2n},K_{2n}\}$ and $\{K_{2n-2},K_{2n+2}\}$.  The first block has the
same Jordan--H\"older length and factors as the known summand $2Q_{2n}$,
so it has no further complement.  In the second block, remove the known
simple summand $Q_{2(n-1)}=K_{2n-2}$; the remaining length-one object is
$K_{2n+2}$.  Hence $Q_{2(n+1)}=K_{2(n+1)}$, completing the induction.

Every injection used in this argument comes either from exact tensoring by
$K_2$ or from the exact projection of Lemma~\ref{lem:block-projection}.  In particular, we
never obtain \eqref{eq:F-projective-filtration} by applying $F$ directly to
\eqref{eq:source-projective-filtration}.
\end{proof}

\begin{rmk}\label{rmk:no-left-exactness}
The final sentence is essential: right exactness alone would only give
right exactness at the last two terms after applying $F$ to
\eqref{eq:source-projective-filtration}.  Proposition
\ref{prop:projective-images} is independent of the exactness ultimately
proved in Theorem~\ref{thm:main}.
\end{rmk}

\section{Logarithmic projective images and full faithfulness}\label{sec:fullness}

Put $Q_j=F(P_j)$.

\begin{lem}\label{lem:p-two-faithful}
If $p=2$, then $F|_{\mathcal P^k}$ is faithful.
\end{lem}

\begin{proof}
Let
\[
 \Phi:\mathcal T_\zeta\xrightarrow{\sim}\mathcal P^k,
 \qquad \Phi(T_\lambda)=P_{\lambda+1},
\]
be the equivalence of Theorem~6.6 and Proposition~6.4 of
\cite{McRaeYang}, now specialized to $p=2$.  In this case the non-wall
projectives form the single chain
\[
 P_1\;\substack{\longrightarrow\\[-3pt]\longleftarrow}\;
 P_3\;\substack{\longrightarrow\\[-3pt]\longleftarrow}\;
 P_5\;\substack{\longrightarrow\\[-3pt]\longleftarrow}\;\cdots,
\]
and the even projectives $P_{2n}$ are isolated simple-projectives; this is
the $p=2$ specialization of the projective Hom algebra in Sections~4 and 6
of \cite{McRaeYang}.  No identity of an indecomposable projective is
killed.  Indeed, Lemma~\ref{lem:F-small-Weyl} gives
$F(P_1)=F(V_1)\cong K_1\neq0$, while
Proposition~\ref{prop:projective-images} gives $F(P_j)=Q_j\neq0$ for every
$j\geq2$.

We next show that neither of the two nonzero arrows between adjacent odd
projectives is killed.  Set $G=F\circ\Phi$.  Under
\eqref{eq:tilting-projective-equivalence}, these arrows are the maps
\[
 f_m^+:T_{2m}\longrightarrow T_{2m+2},
 \quad m\geq0,
 \qquad
 f_m^-:T_{2m}\longrightarrow T_{2m-2},
 \quad m\geq1,
\]
in the notation of Section~6.1 of \cite{McRaeYang}.  The two
rigidity-adjunction calculations in the proof of Theorem~6.6 there,
evaluated at $p=2$, give isomorphisms
\begin{align*}
 \Hom(T_{2m},T_{2m+2})
 &\xrightarrow{\ \sim\ }\End(T_{2m+1}),
 &f_m^+&\longmapsto\id_{T_{2m+1}},\\
 \Hom(T_{2m},T_{2m-2})
 &\xleftarrow{\ \sim\ }\End(T_{2m-1}),
 &f_m^-&\longmapsfrom\id_{T_{2m-1}}.
\end{align*}
The outer isomorphisms use the multiplicity-one summands in
$X\otimes T_{2m}$ and $X\otimes T_{2m\pm1}$, and the middle step is the
adjunction for the rigid self-dual object $X$.  Lemma~6.5 of
\cite{McRaeYang} proves that these adjunctions commute with a tensor
functor preserving the chosen evaluation and coevaluation.  The required
normalization of $G=F\circ\Phi$ has two inputs: the equivalence $\Phi$ is
normalized on the chosen duality morphisms in the construction and proof
of \cite[Theorem~6.6]{McRaeYang}, and $F$ is constructed to carry the chosen
evaluation and coevaluation of $V_2$ to those of $K_{2,1}$ in
\cite[Theorem~7.15]{McRaeYang}.  Hence $G(f_m^+)=0$ would force
$G(\id_{T_{2m+1}})=0$, and $G(f_m^-)=0$ would force
$G(\id_{T_{2m-1}})=0$.  Both are impossible because these identities map
to the identities of the nonzero even-wall objects $Q_{2m+2}$ and
$Q_{2m}$.  Thus all adjacent arrows survive under $F$.

For clarity, the naturality assertion just used is the commutativity, for
each sign, of the square
\[
 \begin{array}{ccc}
  \Hom(T_{2m},T_{2m\pm2})&\xrightarrow{\ \mathcal A_m^\pm\ }&
       \End(T_{2m\pm1})\\[2mm]
  \big\downarrow G_*&&\big\downarrow G_*\\[2mm]
  \Hom(GT_{2m},GT_{2m\pm2})&\xrightarrow{\ \mathcal A_{m,G}^\pm\ }&
       \End(GT_{2m\pm1}),
 \end{array}
\]
where the adjunctions are normalized so that
$\mathcal A_m^\pm(f_m^\pm)=\id_{T_{2m\pm1}}$; here
$\mathcal A_m^-$ denotes the inverse of the second isomorphism displayed
above.  The preceding duality
normalizations are exactly what makes the vertical maps in this square the
maps induced by $G$, rather than unspecified nonzero scalar multiples.

It remains only to check the radical endomorphism of each $P_{2m+1}$,
$m\geq1$.  The source recursion gives
$V_2\boxt P_{2m}\cong P_{2m+1}$.  Applying the monoidal functor $F$ and
using Proposition~\ref{prop:projective-images} for the even-wall object
therefore gives
\[
 Q_{2m+1}\cong K_2\boxt K_{2m}.
\]
Since $2\mid 2m$, McRae--Sopin's wall fusion theorem says that this module
is logarithmic
\cite[Theorem~6.5 and Remark~6.11]{McRaeSopin}.  Preservation of the minus
twist and \eqref{eq:source-twist-nilpotent} give
\[
 \theta_{Q_{2m+1}}
 =\lambda_{2m+1}
  \bigl(\id+\gamma_{2m+1}F(N_{2m+1})\bigr).
\]
The left-hand side is non-semisimple, so $F(N_{2m+1})\neq0$.  Moreover,
\[
 F(N_{2m+1})^2=F(N_{2m+1}^2)=0.
\]
Thus $F(N_{2m+1})$ is a nonzero nilpotent and cannot be a scalar multiple
of $F(\id)=\id_{Q_{2m+1}}$.  Hence these two images are linearly independent,
and the induced map on
\[
 \End(P_{2m+1})=\C\id\oplus\C N_{2m+1}
\]
is injective.  The identities, the adjacent arrows, and these radical loops
form bases of all nonzero Hom spaces in the $p=2$ projective zigzag
category.  Hence no nonzero morphism is killed, proving faithfulness.
\end{proof}

\begin{cor}\label{cor:F-faithful-projectives}
For every $p\geq2$, the restriction of $F$ to the projective subcategory is
faithful.
\end{cor}

\begin{proof}
Combine Proposition~\ref{prop:F-faithful-projectives-large} with
Lemma~\ref{lem:p-two-faithful}.
\end{proof}

\begin{prop}\label{prop:Q-logarithmic}
For $n\geq1$ and $1\leq r\leq p-1$, the module $Q_{np+r}$ is logarithmic.
In particular, the sequence \eqref{eq:F-projective-filtration} is non-split.
\end{prop}

\begin{proof}
By Corollary~\ref{cor:F-faithful-projectives},
$F(N_{np+r})\neq0$.  Since $F$ preserves the minus twist, applying it to
\eqref{eq:source-twist-nilpotent} yields
\[
 \theta_{Q_{np+r}}=
 \lambda_{np+r}(\id+\gamma_{np+r}F(N_{np+r})),
\]
which is not semisimple.  A split sum of the two ordinary Kac modules in
\eqref{eq:F-projective-filtration} would have semisimple $L_0$, so the
sequence cannot split.
\end{proof}

\begin{cor}\label{cor:Q-Ptau}
For $n\geq1$ and $1\leq r\leq p-1$,
\[
 Q_{np+r}\cong P(\tau_{n,r}).
\]
Moreover, if $S_{np+r}\subset K_{np-r}\subset Q_{np+r}$ denotes the
distinguished copy furnished by the Kac submodule, then
\[
 Q_{np+r}/S_{np+r}
\]
is indecomposable.
\end{cor}

\begin{proof}
Let $a,b,c$ be as in \eqref{eq:abc}, and let
$S_b\subset K_a\subset Q_b$ be the copy in the Kac submodule.  The quotient
$Q_b/S_b$ fits into
\[
 0\longrightarrow S_a\longrightarrow Q_b/S_b
 \longrightarrow K_b\longrightarrow0.
\]
Lemma~\ref{lem:Felder-labels} and \eqref{eq:A-B-tau} identify
$K_a=A(\tau_{n,r})$ and $K_b=B(\tau_{n,r})$.  Moreover, all terms and both
short exact sequences lie in the extension-closed category
$\mathscr L_{p_+,p_-}$ fixed before Theorem~\ref{thm:Nakano-input}.
Consequently
\[
 0\longrightarrow A(\tau_{n,r})\longrightarrow Q_b
 \longrightarrow B(\tau_{n,r})\longrightarrow0
\]
represents a class in the precise Yoneda group to which Nakano's
Theorem~5.11 applies, namely
\[
 \Ext^1_{\mathscr L_{p_+,p_-}}
 \bigl(L(h_{\alpha_2},h_{\alpha_3}),
       L(h_{\alpha_1},h_{\alpha_2})\bigr).
\]
By Proposition~\ref{prop:Q-logarithmic}, $L_0$ is
non-semisimple on $Q_b$; this is exactly ``logarithmic'' in Nakano's
Definition~3.14, not an additional Loewy-theoretic condition.
Theorem~\ref{thm:Nakano-input} therefore says that $Q_b/S_b$ is
indecomposable.

The quotient sequence displayed above is thus non-split.  The relevant
simple-by-Kac extension space is one-dimensional by Nakano's
Proposition~5.5, as recorded in the definition of $K(\tau)$ before
Theorem~\ref{thm:Nakano-input}; hence
$Q_b/S_b\cong K(\tau_{n,r})$.  We consequently have
\[
 0\longrightarrow S_b\longrightarrow Q_b
 \longrightarrow K(\tau_{n,r})\longrightarrow0.
\]
This extension is nonzero.  Indeed, in the paragraph preceding
Definition~5.6, Nakano specifies that the indecomposable modules introduced
in Definitions~5.9--5.10, including $K(\tau)$, are quotients of Virasoro
Verma modules \cite[Section~5.2]{Nakano}.  They are therefore ordinary
$L_0$-graded lowest-weight modules; in particular, $L_0$ acts semisimply on
$K(\tau_{n,r})$.  A split middle term
$S_b\oplus K(\tau_{n,r})$ would therefore have semisimple $L_0$, contrary to
Proposition~\ref{prop:Q-logarithmic}.  The one-dimensional
$\Ext^1_{\mathscr L_{p_+,p_-}}$ statement in
Theorem~\ref{thm:Nakano-input}
therefore identifies $Q_b$ with $P(\tau_{n,r})$.  The quotient assertion is
Nakano's Theorem~5.11 applied above.
\end{proof}

We now compute the Hom spaces without using a socle description of
$P(\tau)$.  Fix $1\leq r\leq p-1$ and define a reflection chain
\begin{equation}\label{eq:reflection-chain}
 s_{2m}=2mp+r,
 \qquad
 s_{2m+1}=2(m+1)p-r
 \qquad(m\geq0).
\end{equation}
Write
\[
 \mathsf S_i=S_{s_i},
 \qquad \mathsf A_i=K_{s_i},
 \qquad \mathsf Q_i=Q_{s_i}.
\]

\begin{lem}\label{lem:chain-filtrations}
For every $i\geq0$,
\begin{equation}\label{eq:A-chain}
 0\longrightarrow\mathsf S_{i+1}\longrightarrow\mathsf A_i
 \longrightarrow\mathsf S_i\longrightarrow0
\end{equation}
is non-split.  For every $i\geq1$,
\begin{equation}\label{eq:Q-chain}
 0\longrightarrow\mathsf A_{i-1}\longrightarrow\mathsf Q_i
 \longrightarrow\mathsf A_i\longrightarrow0,
\end{equation}
and $\mathsf Q_0=\mathsf A_0$.
\end{lem}

\begin{proof}
If $s_i=n_ip+r_i$ with $1\leq r_i\leq p-1$, direct inspection of
\eqref{eq:reflection-chain} gives
\[
 s_{i-1}=n_ip-r_i\quad(i\geq1),
 \qquad s_{i+1}=(n_i+2)p-r_i.
\]
Now \eqref{eq:A-chain} is \eqref{eq:Kac-length-two} and
\eqref{eq:Q-chain} is Proposition~\ref{prop:projective-images}.
\end{proof}

\begin{lem}\label{lem:Q-simple-Hom}
For $i\geq1$,
\begin{equation}\label{eq:Q-simple-Hom}
 \Hom(\mathsf Q_i,\mathsf S_j)=
 \begin{cases}
 \C,&j=i,\\
 0,&j\neq i.
 \end{cases}
\end{equation}
\end{lem}

\begin{proof}
For a non-split length-two sequence
$0\to T'\to A\to T\to0$ with $T\not\cong T'$ simple, one has
\begin{equation}\label{eq:length-two-Hom}
 \End(A)=\C,
 \qquad \Hom(A,T')=0,
 \qquad \Hom(A,T)=\C.
\end{equation}
Apply $\Hom(-,\mathsf S_j)$ to \eqref{eq:Q-chain}.  Composition factors and
\eqref{eq:length-two-Hom} immediately give the asserted vanishing except
possibly at $j=i-1$ and $j=i$; the quotient
$\mathsf Q_i\twoheadrightarrow\mathsf A_i\twoheadrightarrow\mathsf S_i$
gives the one-dimensional space at $j=i$.

It remains to prove $\Hom(\mathsf Q_i,\mathsf S_{i-1})=0$.  The relevant
part of the long exact sequence is
\[
 \Hom(\mathsf A_{i-1},\mathsf S_{i-1})
 \xrightarrow{\delta}
 \Ext^1(\mathsf A_i,\mathsf S_{i-1}).
\]
The connecting map sends the quotient
$\mathsf A_{i-1}\twoheadrightarrow\mathsf S_{i-1}$ to the pushout
\begin{equation}\label{eq:pushout}
 0\longrightarrow\mathsf S_{i-1}\longrightarrow
 \mathsf Q_i/\mathsf S_i\longrightarrow\mathsf A_i
 \longrightarrow0.
\end{equation}
By Corollary~\ref{cor:Q-Ptau}, the middle term of
\eqref{eq:pushout} is indecomposable, so the extension is non-split and
$\delta\neq0$.  Its domain is one-dimensional; hence $\delta$ is injective,
which proves the remaining vanishing.
\end{proof}

\begin{lem}\label{lem:Q-Kac-Hom}
For $i\geq1$,
\[
 \Hom(\mathsf Q_i,\mathsf A_j)=
 \begin{cases}
 \C,&j=i-1\text{ or }j=i,\\
 0,&\text{otherwise}.
 \end{cases}
\]
\end{lem}

\begin{proof}
Apply $\Hom(\mathsf Q_i,-)$ to \eqref{eq:A-chain}.  Lemma
\ref{lem:Q-simple-Hom} gives all vanishings.  At $j=i-1$, the injection
$\mathsf S_i\hookrightarrow\mathsf A_{i-1}$ gives the nonzero map.  At
$j=i$, the quotient $\mathsf Q_i\twoheadrightarrow\mathsf A_i$ gives the
nonzero map.  The same exact sequences bound both dimensions by one.
\end{proof}

\begin{prop}\label{prop:target-Hom-table}
Within a fixed reflection chain,
\begin{equation}\label{eq:target-Hom-table}
 \dim\Hom(\mathsf Q_i,\mathsf Q_j)=
 \begin{cases}
 2,&i=j\geq1,\\
 1,&i=j=0,\\
 1,&|i-j|=1,\\
 0,&|i-j|\geq2.
 \end{cases}
\end{equation}
\end{prop}

\begin{proof}
First let $i\geq1$.  Apply $\Hom(\mathsf Q_i,-)$ to
\eqref{eq:Q-chain} for the target.  Lemma~\ref{lem:Q-Kac-Hom} gives the zero
spaces for $|i-j|\geq2$ and gives dimension at most one for adjacent indices.
The source category contains a nonzero morphism in each direction between
adjacent projectives, and its image is nonzero by Corollary
\ref{cor:F-faithful-projectives}; hence the adjacent dimensions equal one.

For $i\geq1$, the same exact sequence gives
\[
 \dim\End(\mathsf Q_i)\leq
 \dim\Hom(\mathsf Q_i,\mathsf A_{i-1})+
 \dim\Hom(\mathsf Q_i,\mathsf A_i)=2.
\]
The identity and the nonzero nilpotent $F(N_{s_i})$ are linearly
independent, so equality holds.

For the boundary source $i=0$, use $\mathsf Q_0=\mathsf A_0$.  If $j\geq3$,
the composition factors are disjoint.  If $j=2$, a map
$\mathsf A_0\to\mathsf Q_2$ has zero composite with
$\mathsf Q_2\twoheadrightarrow\mathsf A_2$, so it lands in
$\mathsf A_1$; but $\Hom(\mathsf A_0,\mathsf A_1)=0$ by
\eqref{eq:A-chain}, \eqref{eq:length-two-Hom}, and composition factors.
Thus $\Hom(\mathsf Q_0,\mathsf Q_j)=0$ for $j\geq2$.  For $j=1$, applying
$\Hom(\mathsf A_0,-)$ to
$0\to\mathsf A_0\to\mathsf Q_1\to\mathsf A_1\to0$ gives dimension at most
one because $\End(\mathsf A_0)=\C$ and
$\Hom(\mathsf A_0,\mathsf A_1)=0$.  The nonzero adjacent source morphism and
faithfulness of $F$ give equality.  Finally, $\mathsf A_0$ is a non-split
length-two module with distinct composition factors, so its endomorphisms
are scalar by \eqref{eq:length-two-Hom}.  Finally, since
$\mathsf Q_0=\mathsf A_0$, the cases with target index $0$ and positive
source index follow directly from Lemma~\ref{lem:Q-Kac-Hom}.
This proves every boundary entry.
\end{proof}

We must also rule out morphisms between distinct reflection chains.  For
$q>2$, this follows directly from composition factors.  The only collision
requiring attention occurs when $q=2$.

\begin{lem}\label{lem:cross-chain}
The Hom spaces between projective images belonging to distinct reflection
chains are zero.
\end{lem}

\begin{proof}
From \eqref{eq:weights},
\begin{equation}\label{eq:weight-equality}
 h_j=h_\ell
 \quad\Longrightarrow\quad
 j=\ell\quad\text{or}\quad q(j+\ell)=2p.
\end{equation}
Because $(p,q)=1$ and $q\geq2$, the second alternative occurs only for
$q=2$, when $p$ is odd and $h_r=h_{p-r}$.  Thus for $q>2$ distinct chains
have disjoint simple composition factors.

Assume $q=2$ and compare the chains beginning at $r$ and $p-r$, with
$1\leq r<p-r\leq p-1$.  Their only common simple is
$S:=S_r\cong S_{p-r}$.  Put $A=K_r$ and $B=K_{p-r}$.  Their non-split
sequences are
\[
 0\to S_{2p-r}\to A\to S\to0,
 \qquad
 0\to S_{p+r}\to B\to S\to0,
\]
with nonisomorphic submodules.  Any map $A\to B$ composed with $B\to S$
is a scalar multiple of $A\to S$.  If this scalar is zero, the map lands in
$S_{p+r}$ and is zero by composition factors.  If it is nonzero, the map
kills $S_{2p-r}$ and factors through a section $S\to B$, contradicting
nonsplitting.  Hence $\Hom(A,B)=0$; similarly $\Hom(B,A)=0$.

Let $Q^r_i$ and $Q^{p-r}_j$ denote the two chains.  The only further pairs
that share a composition factor are those for which one index is $0$ or
$1$.  We show, for example, that
$\Hom(Q^r_1,B)=\Hom(B,Q^r_1)=0$.  By
Lemma~\ref{lem:Q-simple-Hom}, every map $Q^r_1\to B$ has zero composite
with $B\twoheadrightarrow S$; it therefore lands in the simple submodule
$S_{p+r}$, which is not a composition factor of $Q^r_1$.  Conversely,
compose a map $B\to Q^r_1$ with the quotient
$Q^r_1\twoheadrightarrow K_{2p-r}$.  The composition is zero by disjoint
composition factors, so the map lands in $A\subset Q^r_1$ and vanishes by
$\Hom(B,A)=0$.  The analogous assertions with $r$ and $p-r$ exchanged also
hold.

For the pair $(Q^r_1,Q^{p-r}_1)$, apply
$\Hom(Q^r_1,-)$ to the target filtration
$0\to B\to Q^{p-r}_1\to K_{p+r}\to0$.  The Hom space to $B$ vanishes by
the preceding paragraph, and the Hom space to $K_{p+r}$ vanishes after
applying Lemma~\ref{lem:Q-simple-Hom} to its two simple factors.  Thus
$\Hom(Q^r_1,Q^{p-r}_1)=0$, and the reverse direction is identical.  For all
remaining pairs the composition-factor sets are disjoint.  This proves all
cross-chain vanishings.
\end{proof}

\begin{lem}[Wall objects]\label{lem:wall-Hom}
For $m,n\geq1$,
\[
 \Hom(Q_{np},Q_{mp})=
 \begin{cases}
  \C,&m=n,\\
  0,&m\neq n,
 \end{cases}
\]
and every Hom space between a wall object $Q_{np}$ and a non-wall
projective image is zero, in either direction.
\end{lem}

\begin{proof}
Proposition~\ref{prop:projective-images} gives
$Q_{np}=K_{np}=S_{np}$, so a wall object is simple.  Every composition
factor of a non-wall $Q_j$ has label not divisible by $p$, by
\eqref{eq:source-composition-factors} and the projective-image filtration.
Here a composition factor is indexed by its simple-module isomorphism
class, so possible equality of conformal weights must still be excluded.
Formula \eqref{eq:weight-equality} says that a wall label $np$ and a
non-wall label $j$ could collide only if
\[
 q(j+np)=2p.
\]
Since $j,np>0$, this forces $n=1$, $q=2$, and $j=0$, which is not a
non-wall label.  Thus the simple composition-factor sets are disjoint,
proving both directional vanishings.  For two wall labels, the same second
alternative would say $q(n+m)=2$; this is impossible for $q\geq2$ and
$n,m\geq1$ unless $q=2$ and $n=m=\tfrac12$.  Hence distinct wall objects
are nonisomorphic.  Their simplicity completes the proof.
\end{proof}

\begin{thm}\label{thm:F-fully-faithful}
The restriction
\[
 F|_{\mathcal P^k}:\mathcal P^k\longrightarrow\Oc_{c_{p,q}}
\]
to the projective subcategory is fully faithful.
\end{thm}

\begin{proof}
McRae--Yang's projective Hom algebra has the zigzag dimensions in
\eqref{eq:target-Hom-table}: scalar endomorphisms at the boundary,
two-dimensional endomorphisms at every nontrivial projective, one-dimensional
Hom spaces in both directions between adjacent projectives, and zero Hom
spaces otherwise; see \cite[Sections~4 and 6]{McRaeYang}.  Simple-projective
    objects $P_{np}$ are isolated and have scalar endomorphisms.  Proposition
    \ref{prop:target-Hom-table} and Lemmas~\ref{lem:cross-chain} and
    \ref{lem:wall-Hom} show that the target Hom spaces have the same
    dimensions, including all boundary and wall cases.  Corollary
\ref{cor:F-faithful-projectives} makes each induced map on Hom spaces
injective, so each is an isomorphism.  Finite direct sums follow by matrices.
\end{proof}

\section{Comparison with Drinfeld--Sokolov reduction}\label{sec:comparison}

We first compare objects.  Exactness of $H$ and
Proposition~\ref{prop:H-Weyl} applied to
\eqref{eq:source-projective-filtration} give
\begin{equation}\label{eq:H-projective-filtration}
 0\longrightarrow K_{np-r}\longrightarrow H(P_{np+r})
 \longrightarrow K_{np+r}\longrightarrow0.
\end{equation}

\begin{prop}\label{prop:projective-object-comparison}
For every indecomposable projective $P_j$,
\[
 F(P_j)\cong H(P_j).
\]
\end{prop}

\begin{proof}
If $1\leq j\leq p-1$, then $P_j=V_j$, and
\[
 F(P_j)=F(V_j)\cong K_j\cong H(V_j)=H(P_j)
\]
by Lemma~\ref{lem:F-small-Weyl} and Proposition~\ref{prop:H-Weyl}.
If $p\mid j$, both sides are $K_j$ by
Propositions~\ref{prop:projective-images} and \ref{prop:H-Weyl}.
Let $j=np+r$ with $n\geq1$ and $1\leq r\leq p-1$.  Corollary
\ref{cor:H-faithful} and Proposition~\ref{prop:H-twist}, applied to the
nonzero nilpotent in \eqref{eq:source-twist-nilpotent}, show that the twist
on $H(P_j)$ is non-semisimple.  Hence \eqref{eq:H-projective-filtration} is
a logarithmic extension.  Repeating the quotient argument in the proof of
Corollary~\ref{cor:Q-Ptau}, Theorem~\ref{thm:Nakano-input} first identifies
$H(P_j)/S_j$ with $K(\tau_{n,r})$ and then uses the one-dimensional
extension space to give
\[
 H(P_j)\cong P(\tau_{n,r})\cong F(P_j),
\]
the second isomorphism being Corollary~\ref{cor:Q-Ptau}.
\end{proof}

Objectwise isomorphisms do not by themselves give naturality.  We use the
nilpotent part of the twist to lock the remaining scalars.

\begin{lem}[Twist locking]\label{lem:twist-locking}
Let $P=P_{np+r}$ be a nontrivial indecomposable projective and let
$\nu_P$ span $\rad\End(P)$.  Every Virasoro-module isomorphism
$\eta_P:F(P)\to H(P)$ satisfies
\[
 \eta_PF(\nu_P)=H(\nu_P)\eta_P.
\]
\end{lem}

\begin{proof}
The twist of $P$ has the form
\[
 \theta^-_P=\lambda_P(\id+\gamma_P\nu_P),
 \qquad \gamma_P\neq0,
\]
by \eqref{eq:source-twist-nilpotent}.  A Virasoro-module isomorphism
commutes with the target twist.  Since $F$ and $H$ both preserve the source
minus twist, by Theorem~7.15 of \cite{McRaeYang} and
Proposition~\ref{prop:H-twist}, respectively,
\[
 \eta_PF(\theta^-_P)=H(\theta^-_P)\eta_P.
\]
Canceling the scalar and identity terms gives the claimed equality.
\end{proof}

\begin{prop}\label{prop:natural-on-projectives}
There is a natural isomorphism
\[
 \eta:F|_{\mathcal P^k}\xRightarrow{\ \sim\ }
 H|_{\mathcal P^k}.
\]
\end{prop}

\begin{proof}
Work first on one reflection chain and write its projectives as
$P_0,P_1,\ldots$.  Choose nonzero maps
\[
 P_i\underset{y_i}{\stackrel{x_i}{\rightleftarrows}}P_{i+1}
\]
so that $\nu_{i+1}=x_iy_i$ is the nonzero radical endomorphism of
$P_{i+1}$; this is the standard projective zigzag normalization in
\cite[Sections~4 and 6]{McRaeYang}.  Both $F$ and $H$ send the four maps
$x_i,y_i$ to nonzero maps: for $F$ this follows from Corollary
\ref{cor:F-faithful-projectives}, and for $H$ from Corollary
\ref{cor:H-faithful}.

Choose $\eta_0:F(P_0)\to H(P_0)$.  Inductively, after $\eta_i$ is chosen,
choose an objectwise isomorphism at $P_{i+1}$ and rescale it so that
\begin{equation}\label{eq:forward-naturality}
 \eta_{i+1}F(x_i)=H(x_i)\eta_i.
\end{equation}
This is possible because the relevant adjacent Hom space is
one-dimensional.  The reverse Hom space is also one-dimensional, so for
some $c_i\in\C^\times$,
\begin{equation}\label{eq:reverse-scalar}
 \eta_iF(y_i)=c_iH(y_i)\eta_{i+1}.
\end{equation}
Composing \eqref{eq:forward-naturality} and
\eqref{eq:reverse-scalar} gives
\[
 \eta_{i+1}F(\nu_{i+1})
 =c_iH(\nu_{i+1})\eta_{i+1}.
\]
Lemma~\ref{lem:twist-locking} gives the same equality with $c_i=1$.
Since $H$ is faithful and $\nu_{i+1}\neq0$, it follows that $c_i=1$.
Thus naturality holds for both arrows.

The identities, adjacent arrows, and radical loops span every nonzero Hom
space in the projective zigzag category, so the resulting family is natural
on the chain.  On isolated simple-projectives, any objectwise isomorphism is
automatically natural.  Distinct chains have no morphisms, and finite direct
sums are handled by matrices.  This proves the proposition.
\end{proof}

\begin{proof}[Proof of Theorem~\ref{thm:main}]
The functor $F$ is right exact by construction, and $H$ is exact, hence
right exact.  The source category has enough projectives.  McRae--Yang's
right-exact extension theorem \cite[Theorem~6.7]{McRaeYang} says that a
$\C$-linear functor on the projective subcategory has a unique, up to natural
isomorphism, right exact extension to the entire abelian category.
Proposition~\ref{prop:natural-on-projectives} therefore extends uniquely to
\[
 F\cong H
\]
on $\KL^k(\mathfrak{sl}_2)$.  Since $H$ is exact, $F$ is exact.

Every result used in the preceding paragraph was proved above for an
arbitrary ordered coprime pair.  Apply that already proved result to the
ordered pair $(q,p)$.  Its level is $k'=-2+q/p$, its McRae--Yang functor is
$F_{q,p}$, and its target generator is
$\Kac_{j,1}^{(q,p)}$.  With superscripts indicating the ordered
central-charge parameters, the Feigin--Fuchs weight formula gives
\[
 h_{j,1}^{(q,p)}=h_{1,j}^{(p,q)}.
\]
This equality upgrades to a module isomorphism, not merely an equality of
lowest weights.  By the definition of Virasoro Kac modules
\cite[Remark~2.2]{McRaeSopin}, both $\Kac_{j,1}^{(q,p)}$ and
$\Kac_{1,j}^{(p,q)}$ are quotients of the same Virasoro Verma module
\[
 \mathcal V\bigl(c_{p,q},h_{j,1}^{(q,p)}\bigr)
 =\mathcal V\bigl(c_{p,q},h_{1,j}^{(p,q)}\bigr)
\]
by the submodule generated by its singular vector at level $j$.  That
singular vector is unique up to scalar
\cite[Proposition~2.2.1(3)]{CreutzigEtAl}, so the two generated submodules
coincide.  Therefore
\[
 \Kac_{j,1}^{(q,p)}\cong\Kac_{1,j}^{(p,q)}.
\]
Since $c_{q,p}=c_{p,q}$, this identifies the target category with the one in
the statement of the theorem.  In particular, when $q=2$ this is a direct
application of Lemma~\ref{lem:p-two-faithful} to the ordered pair whose first
entry is two; no unproved boundary-case transposition is being invoked.
The conclusion already established for an arbitrary ordered pair now gives
\[
 F_{q,p}\cong H^{0\prime}_{DS,+}
 \]
and exactness of $F_{q,p}$.  Together with the first half, this is precisely
\cite[Conjecture~7.16]{McRaeYang}.
\end{proof}

\begin{rmk}
The conclusion is a natural isomorphism of $\C$-linear functors, exactly as
stated in Conjecture~7.16.  Transporting the braided tensor structure of
$F_{p,q}$ along this isomorphism equips the restricted reduction with such a
structure.  Comparing it with any independently constructed monoidal
structure on Drinfeld--Sokolov reduction is a separate question and is not
needed here.
\end{rmk}

\section*{Acknowledgments}

We gratefully acknowledge the use of OpenAI's GPT5.6 Sol model as a research assistance tool in the preparation of this manuscript. The model was employed to support literature exploration, conceptual brainstorming, and language polishing, which contributed to improving the clarity and efficiency of our work.

\begin{thebibliography}{99}

\bibitem{ArakawaVanishing}
T.~Arakawa,
\emph{Vanishing of cohomology associated to quantized Drinfeld--Sokolov reduction},
Int. Math. Res. Not. 2004 (2004), no.~15, 729--767.

\bibitem{ArakawaW}
T.~Arakawa,
\emph{Representation theory of $W$-algebras},
Invent. Math. 169 (2007), no.~2, 219--320.

\bibitem{ArakawaSuperconformal}
T.~Arakawa,
\emph{Representation theory of superconformal algebras and the
Kac--Roan--Wakimoto conjecture},
Duke Math. J. 130 (2005), no.~3, 435--478.

\bibitem{ArakawaFrenkel}
T.~Arakawa and E.~Frenkel,
\emph{Quantum Langlands duality of representations of $\mathcal W$-algebras},
Compos. Math. 155 (2019), no.~12, 2235--2262.

\bibitem{FKW}
E.~Frenkel, V.~Kac and M.~Wakimoto,
\emph{Characters and fusion rules for $W$-algebras via quantized Drinfeld--Sokolov reduction},
Comm. Math. Phys. 147 (1992), no.~2, 295--328.

\bibitem{CreutzigEtAl}
T.~Creutzig, C.~Jiang, F.~Orosz Hunziker, D.~Ridout and J.~Yang,
\emph{Tensor categories arising from the Virasoro algebra},
Adv. Math. 380 (2021), Paper No.~107601, 34 pp.

\bibitem{GoodmanWenzl}
F.~M.~Goodman and H.~Wenzl,
\emph{Ideals in the Temperley--Lieb category},
Appendix to M.~H.~Freedman, \emph{A magnetic model with a possible Chern--Simons phase},
Comm. Math. Phys. 234 (2003), no.~1, 129--183.

\bibitem{McRaeSopin}
R.~McRae and V.~Sopin,
\emph{Fusion and (non)-rigidity of Virasoro Kac modules in logarithmic minimal models at $(p,q)$-central charge},
Phys. Scr. 99 (2024), 035233, 42 pp.

\bibitem{McRaeYang}
R.~McRae and J.~Yang,
\emph{The non-semisimple Kazhdan--Lusztig category for affine $\mathfrak{sl}_2$ at admissible levels},
Proc. Lond. Math. Soc. (3) 130 (2025), e70043.

\bibitem{Nakano}
H.~Nakano,
\emph{Projective covers of the simple modules for the triplet $W$-algebra $\mathcal W_{p_+,p_-}$},
arXiv:2305.12448v2 [math.RT], 9 August 2026.

\bibitem{SuttonTubbenhauerWedrichZhu}
L.~Sutton, D.~Tubbenhauer, P.~Wedrich and J.~Zhu,
\emph{$\mathrm{SL}_2$ tilting modules in the mixed case},
Selecta Math. (N.S.) 29 (2023), no.~3, Paper No.~39, 40 pp.

\end{thebibliography}

\end{document}
